% answers_short.tex —— 1991 年·上海卷 · 简短答案速查
% 由 tools/make_short_answers.py 自 answers.tex 的 \gksolhead/\gkanswer 提取，
% 并与 1991/1991_上海_答案要点.md 逐题比对后生成，勿手改（改源 MD 后重跑）。

\begin{gknotice}
\textbf{使用说明}：本节只给最终答案便于快速核对，完整推导见解析版。答案为编写者独立推导，\textbf{非官方标准答案}。
\end{gknotice}

\gkpart{填空题}
\begin{itemize}[leftmargin=2.6em]
  \item[\textbf{1.}] $y=\log_{3}x$（$x\ge 1$）。
  \item[\textbf{2.}] $-\dfrac{\sqrt{10}}{10}$。
  \item[\textbf{3.}] $\dfrac{3}{2}$。
  \item[\textbf{4.}] $2$。
  \item[\textbf{5.}] $\dfrac{\sqrt{3}}{2}$。
  \item[\textbf{6.}] $\dfrac{9}{2}\sqrt{3}$。
  \item[\textbf{7.}] $8\sqrt{2}\,\pi$。
  \item[\textbf{8.}] 无常数项（题面指数 $8$；卷内原文作 $60$，对应指数 $6$，源疑点见 ISSUE-1178）。
  \item[\textbf{9.}] $y=x^{2}-4x+5$。
  \item[\textbf{10.}] $\rho=-8\cos\theta$。
\end{itemize}

\gkpart{选择题}
\begin{center}\begin{tabular}{ccccccccccc}
\toprule
题号 & 11 & 12 & 13 & 14 & 15 & 16 & 17 & 18 & 19 & 20 \\
\midrule
答案 & B & C & B & C & B & A & D & A & D & A \\
\bottomrule
\end{tabular}\end{center}

\gkpart{解答题}
\begin{itemize}[leftmargin=2.6em]
  \item[\textbf{21.}] $\cos 4x=\dfrac{1}{2}$。
  \item[\textbf{22.}] $\dfrac{2-\sqrt{3}}{4}d$ 或 $\dfrac{2+\sqrt{3}}{4}d$（万公里）。
  \item[\textbf{23.}] 当 $0<a<\dfrac{1}{2}$ 时，$x>0$；当 $a=\dfrac{1}{2}$ 时，$0<x<\dfrac{1}{2}$ 或 $\dfrac{1}{2}<x<2$ 或 $x>2$；当 $\dfrac{1}{2}<a<1$ 时，$0<x<a^{\sqrt{2a}}$ 或 $a^{\sqrt{2(1-a)}}<x<a^{-\sqrt{2(1-a)}}$ 或 $x>a^{-\sqrt{2a}}$；当 $a>1$ 时，$0<x<a^{-\sqrt{2a}}$ 或 $x>a^{\sqrt{2a}}$。
  \item[\textbf{24.}] （1）$45^{\circ}$；（2）$90^{\circ}$；（3）$\pi-\operatorname{arctg}2$。
  \item[\textbf{25.}] （1）$\left(3\theta-\dfrac{\pi}{2}\right)+2k\pi$（$k\in Z$）；（2）以 $(0,1)$ 为圆心、1 为半径的圆，除去点 $(0,2)$。
  \item[\textbf{26.}] 轨迹为直线 $12X-3Y-4=0$ 介于 $4-\dfrac{4}{3}\sqrt{6}<Y<4+\dfrac{4}{3}\sqrt{6}$ 间的一段，且除去点 $\left(\dfrac{4}{3},4\right)$。
  \item[\textbf{27.}] （1）$T=\dfrac{ab}{(a+b)c}$；（2）$T=\dfrac{1}{2(1+\cos\theta)}\left(x-\dfrac{1}{x}\right)$，定义域为 $\left(1,\sqrt{\dfrac{2}{1-\cos\theta}}\right]$，最大值 $u=\dfrac{\sqrt{2}}{4\sqrt{1-\cos\theta}}$；（3）$u$ 的最大值为 $\dfrac{1}{2}$。
\end{itemize}
